% =========================================================
% mechlib/core/styles.tex
% ZENTRALE, NAMENSRAUMGESCHUETZTE TIKZ-STILE
% =========================================================
% Diese Datei definiert ausschliesslich visuelle Grundstile. Die eigentliche
% Geometrie eines Elements gehoert in elements/ bzw. annotations/.
%
% Alle oeffentlichen Bibliotheksstile beginnen mit "mech/". Das reduziert
% Namenskollisionen mit dem Hauptdokument und macht im TikZ-Code sichtbar,
% welche Formatierung aus der mechlib stammt.
%
% Designregel:
%   - Farben werden aus core/colors.tex bezogen.
%   - Geometrien werden NICHT hier definiert.
%   - Linienstarken werden hier zentral gehalten, soweit sie elementuebergreifend
%     konsistent sein sollen.
% =========================================================
% -------------------------------------------------------------------------
% Visuelle Basisstile
% -------------------------------------------------------------------------
% Einzelne Pics duerfen diese Stile ueber ihre jeweiligen "style"-Keys ersetzen.
\tikzset{
  mech/body/.style={
    draw=black,
    fill=mechBodyColor,
    line width=.6pt
  },
  mech/support/.style={
    draw=black,
    fill=mechSupportColor,
    line width=.6pt
  },
  mech/connector/.style={
    draw=black,
    line width=.6pt
  },
  mech/rod/.style={
    draw=black,
    line width=.85pt
  },
  mech/load/.style={
    draw=mechLoadColor,
    text=mechLoadColor,
    line width=1pt
  },
  mech/load fill/.style={
    fill=mechLoadColor,
    fill opacity=.16,
    draw=none
  },
  mech/dimension/.style={
    {Latex[length=2mm]}-{Latex[length=2mm]},
    line width=.5pt
  },
  mech/axis/.style={
    draw=mechAxisColor,
    text=mechAxisColor,
    line width=.6pt
  },
  mech/label/.style={
    inner sep=1pt
  }
}
